\documentclass[a4paper,11pt]{article}

\usepackage[margin=2.3cm]{geometry}
\usepackage[utf8]{inputenc}
\usepackage[T1]{fontenc}
\usepackage{lmodern}
\usepackage{microtype}

\usepackage{graphicx}
\usepackage{subcaption}
\usepackage{amsmath,amssymb}
\usepackage{booktabs}
\usepackage{hyperref}
\usepackage{url}
\usepackage{xcolor}
\usepackage{listings}
\usepackage[numbers,sort&compress]{natbib}

\hypersetup{
    colorlinks=true,
    linkcolor=black,
    citecolor=blue!60!black,
    urlcolor=blue!60!black,
}

\lstset{
    basicstyle=\footnotesize\ttfamily,
    breaklines=true,
    frame=single,
    numbers=left,
    numberstyle=\tiny\color{gray},
    keywordstyle=\color{blue!60!black},
    commentstyle=\color{green!50!black},
    stringstyle=\color{orange!80!black},
    showstringspaces=false,
}

\title{A Hopfield-based Replay Buffer for DQN:\\
An Experimental Study on CartPole-v1}

\author{\.{I}brahim Emir Tosun \\ {\small Erasmus Student} \\ {\small Bio-Inspired Machine Learning (APBIO), Universidad de Vigo}}

\date{Academic Year 2025--2026}

\begin{document}
\maketitle

\begin{abstract}
\noindent
This report describes a hybrid reinforcement learning agent that couples a
Modern Hopfield Network with Deep Q-Learning (DQN). The Hopfield network
replaces the standard uniform-random replay buffer: transitions are sampled
according to attention weights conditioned on the most recently observed
state, so that past experiences most similar to the current situation are
preferentially replayed. The system is tested on the CartPole-v1 environment.
An initial experiment with the default retrieval temperature $\beta=5$ shows
that the hybrid underperforms the plain DQN baseline by about 56\% in final
return. A hyperparameter sweep then reveals that $\beta$ is by far the most
influential parameter: lowering it to $\beta=1.0$ flips the result, and the
best hybrid configuration slightly outperforms the baseline ($263.7 \pm 82.0$
vs.\ $236.4 \pm 98.6$, $+11.5\%$). The improvement is not statistically
significant ($p=0.68$, Welch's $t$-test), which suggests that on a
well-posed environment such as CartPole the Hopfield buffer offers only a
marginal benefit over random replay. The project's main contribution is
showing, with paired 5-seed experiments and three hyperparameter sweeps,
that a biologically-inspired replay mechanism can match but does not clearly
exceed random replay in this setting, and that its behaviour is strongly
governed by the retrieval temperature.
\end{abstract}

\section{Introduction}
\label{sec:intro}

Deep Q-Networks (DQN) store observed transitions in an experience replay
buffer and, at every training step, draw a mini-batch uniformly at random
from that buffer~\citep{mnih2015human}. Random sampling decorrelates the
transitions seen by the optimiser, which stabilises gradient estimates, but
it ignores any notion of \emph{relevance}: a transition recorded in a
completely unrelated state is equally likely to be drawn as one that
closely resembles the agent's current situation.

A natural question is whether the replay buffer can be made
\emph{content-addressable}. Modern Hopfield Networks, recently reformulated
by \citet{ramsauer2020hopfield}, provide exactly such a mechanism: given a
query vector, they return a softmax-weighted combination of the stored
patterns, with sharpness controlled by an inverse-temperature parameter
$\beta$. Functionally they are identical to attention heads, but framed as
an associative memory with a clear energy interpretation.

This project asks a simple empirical question: if we replace DQN's uniform
replay with a Hopfield-based replay buffer that retrieves transitions
similar to the most recent observation, does the agent learn faster, more
robustly, or to a higher final return? The four families of bio-inspired
methods covered in the course (evolutionary, swarm, immune-inspired,
neuro-inspired) all provide candidate components; I picked a neuro-inspired
one because its mathematics matches modern attention mechanisms closely
and integrates cleanly with a standard deep RL pipeline.

\paragraph{Contributions.}
This report contributes (i) a concrete implementation of
\texttt{HopfieldReplayBuffer} compatible with Stable-Baselines3's DQN;
(ii) a paired 5-seed comparison of hybrid and baseline on CartPole-v1;
(iii) three hyperparameter sweeps (one bio-inspired, one structural, one
RL) that identify the retrieval temperature $\beta$ as the only parameter
with a large effect on performance; and (iv) a follow-up experiment with
the best-found configuration that closes and slightly reverses the initial
gap.

\paragraph{Outline.}
Section~\ref{sec:background} reviews DQN replay buffers and Modern
Hopfield Networks. Section~\ref{sec:methodology} describes the hybrid
architecture and the integration with Stable-Baselines3.
Section~\ref{sec:results} reports the main experiment, the sweeps, and the
best-configuration run. Section~\ref{sec:discussion} addresses the ten
evaluation questions from the project specification.
Section~\ref{sec:conclusion} concludes.

\section{Background}
\label{sec:background}

\subsection{Deep Q-Learning and experience replay}

DQN learns a parametric action-value function $Q_\theta(s,a)$ by minimising
the temporal-difference error
\begin{equation}
    \mathcal{L}(\theta) =
    \mathbb{E}_{(s,a,r,s') \sim \mathcal{B}}
    \left[
        \bigl(r + \gamma \max_{a'} Q_{\bar\theta}(s',a') - Q_\theta(s,a)\bigr)^2
    \right],
\end{equation}
where $\mathcal{B}$ is a replay buffer of past transitions and
$\bar\theta$ are the (slowly updated) target-network parameters. Sampling
from $\mathcal{B}$ at random breaks the temporal correlation between
consecutive states, which is what makes function approximation with a
neural network stable. Prioritised Experience
Replay~\citep{schaul2016prioritized} biases sampling by TD-error and was
an early demonstration that non-uniform replay can help; our
Hopfield-based scheme is a different kind of prioritisation, biased by
state similarity rather than by error.

\subsection{Modern Hopfield Networks}

A classical Hopfield network stores $N$ binary patterns
$\{\xi_i\}_{i=1}^{N}$ and retrieves one of them given a noisy cue by
iterative energy minimisation. Its storage capacity is linear in the
pattern dimension, which severely limits its use as a general-purpose
memory. \citet{ramsauer2020hopfield} introduced a continuous variant in
which retrieval is a single matrix operation
\begin{equation}
    \mathrm{retrieve}(q) = X^\top \,\mathrm{softmax}\bigl(\beta X q\bigr),
    \label{eq:hopfield}
\end{equation}
where $X \in \mathbb{R}^{N \times d}$ is the matrix of stored patterns
(one per row), $q \in \mathbb{R}^d$ is the query, and $\beta > 0$ is an
inverse temperature. This formula is identical to a standard attention
layer with $X$ acting as both keys and values. The storage capacity grows
exponentially with $d$, and the associated energy function is
\begin{equation}
    E(q) = -\tfrac{1}{\beta}\log \textstyle\sum_i \exp(\beta x_i^\top q)
        + \tfrac{1}{2}\|q\|^2,
\end{equation}
which lets us monitor how ``confident'' the memory is about a given query.

\subsection{Why combine them?}

Uniform replay is a very cheap form of prioritisation: it weights every
past transition equally. A Hopfield memory, at the opposite extreme, can
concentrate all weight on a single stored pattern when $\beta$ is large.
Intermediate values of $\beta$ interpolate between the two. The working
hypothesis of this project is that a well-tuned Hopfield replay buffer
lies in a sweet spot between ``fully random'' and ``fully content-addressed''
and should therefore at least match, and perhaps slightly improve on,
plain DQN on a simple environment.

There is a second, subtler motivation. In DQN the target
$y = r + \gamma \max_{a'} Q_{\bar\theta}(s',a')$ is itself a bootstrapped
estimate, and errors in it propagate through the Bellman update. If the
mini-batch is dominated by transitions whose states are close to the
current one, the regression becomes locally specialised and the network
may fit the short-term Q-landscape well at the cost of its global
accuracy. Random sampling averages over the whole replay buffer and
implicitly regularises against this local overfitting. A Hopfield
retrieval with a moderate $\beta$ is therefore only expected to help
when ``relevance'' is more informative than ``coverage'' --- when the
problem has clearly separated regions or rare-but-important transitions.
CartPole-v1 has neither property: it is a low-dimensional, dense-reward
task in which almost every transition carries useful information. This
is a deliberate choice on my part: if the hybrid fails even on an
environment this benign, it is a clear signal that the design is not a
general improvement over random replay.

\section{Methodology}
\label{sec:methodology}

\subsection{Environment and base RL algorithm}

I use the \texttt{CartPole-v1} environment from
Gymnasium~\citep{brockman2016openai}: a 4-dimensional continuous state,
two discrete actions (push left or right), reward of $+1$ per step, and a
maximum return of $500$. The base RL algorithm is DQN as implemented in
Stable-Baselines3 version 2.3~\citep{raffin2021stable}. The hyperparameters
used throughout are the Stable-Baselines3 RL Zoo defaults for
CartPole-v1~\citep{rlzoo2020}: learning rate $2.3\!\times\!10^{-3}$,
two hidden layers of $256$ units, $\mathrm{train\_freq}=256$ steps,
$\mathrm{gradient\_steps}=128$, target update every $10$ training steps,
exploration fraction $0.16$, and a buffer of $100\,000$ transitions. I
train for $50\,000$ environment steps per run, which is enough for the
baseline to reach a final return of several hundred on most seeds but not
enough to always saturate at $500$; this leaves headroom for the hybrid
to improve on the baseline.

\subsection{Hopfield replay buffer}

The Hopfield memory is a small PyTorch module with a pattern matrix
$X \in \mathbb{R}^{N \times d}$, an inverse temperature $\beta$, and three
operations:

\begin{itemize}
    \item \texttt{store(x)}: overwrite the oldest patterns if the memory is
    full (ring-buffer); otherwise append.
    \item \texttt{retrieve\_weights(q)}: return
    $\mathrm{softmax}(\beta X q)$.
    \item \texttt{mean\_energy(q)}: return the energy associated with the
    query, used as a diagnostic.
\end{itemize}

\texttt{HopfieldReplayBuffer} subclasses
\texttt{stable\_baselines3.common.buffers.ReplayBuffer}. Its contract is:

\begin{enumerate}
    \item On \texttt{add()}, the transition is appended to the parent
    class's arrays \emph{and} the observation is appended to the Hopfield
    memory. The most recent observation is cached in \texttt{self.\_last\_obs}.
    \item On \texttt{sample(batch\_size)}, the buffer computes the
    attention weights of \texttt{self.\_last\_obs} over all currently
    stored patterns, uses those weights to draw \texttt{batch\_size}
    indices with replacement, and delegates the actual tensor
    construction to the parent class's \texttt{\_get\_samples()}.
\end{enumerate}

\noindent
Pseudocode for the sampling step:

\begin{lstlisting}[language=Python]
def sample(self, batch_size, env=None):
    w, n = self.hopfield.retrieve_weights(self._last_obs)
    w = w.squeeze(0).cpu().numpy()
    w = w[:min(n, self.size())]
    w = w / w.sum()
    idx = np.random.choice(len(w), size=batch_size, replace=True, p=w)
    return self._get_samples(idx, env=env)
\end{lstlisting}

\subsection{Architecture overview}

Figure~\ref{fig:architecture} shows how the components interact. The
environment produces transitions in the usual way, and the DQN policy is
unchanged. The only difference is that the bridge between the
environment's transition stream and the DQN training loop goes through
the Hopfield memory rather than a flat array.

\begin{figure}[t]
    \centering
    \includegraphics[width=0.92\linewidth]{figures/architecture.png}
    \caption{Architecture of the hybrid system. Solid arrows denote the
    normal DQN data flow; dashed arrows highlight the changes introduced
    by the Hopfield replay buffer. The only point of integration is the
    replay buffer itself.}
    \label{fig:architecture}
\end{figure}

\subsection{Validity of the coupling}

The project specification gives three validity checks for a hybrid design:

\begin{description}
    \item[Is the coupling non-trivial?] Yes. Removing the Hopfield
    component and falling back to the parent class's random sampling
    recovers plain DQN exactly; the two systems differ in what
    \texttt{sample()} returns.
    \item[Does the bio-inspired component have adaptive state?] Yes.
    The pattern matrix $X$ grows and is rewritten as new transitions
    arrive, so the set of retrievable patterns evolves continuously
    during training.
    \item[Can the data flow be drawn?] Yes, see
    Figure~\ref{fig:architecture}. The bridge is localised to a single
    module; no other part of DQN is modified.
\end{description}

\subsection{Experimental protocol}

All reported numbers are averages over $5$ independent seeds for the
main experiment and $3$ seeds per configuration for the sweeps. Each run
is a full $50\,000$-step training on a fresh environment and a freshly
initialised network, with deterministic seeds set in Python, NumPy,
PyTorch, and the Gymnasium environment. Raw per-episode CSV logs are
written to disk as each run completes, so experiments are restartable.

The 5-seed budget is reserved for the headline claims (main experiment
and best-configuration validation), in line with the project
specification's requirement that ``all reported results'' must be averaged
over at least 5 seeds. The hyperparameter sweeps use 3 seeds because
their role is diagnostic: they identify which knob has a large effect
and roughly where its sweet spot lies, not to make a publication-grade
claim about any specific configuration. The trends they reveal (most
notably that $\beta$ is the dominant hyperparameter) are then validated
by re-running the best-found configuration with the full 5-seed budget,
which is the result reported in Section~\ref{sec:results}.

\subsection{Departures from the worked example}

The project specification offers ``Hopfield replay buffer + DQN'' as a
\emph{worked example}. My system uses the same insertion point but
differs in several deliberate ways:

\begin{itemize}
    \item \textbf{Sampling rule.} The worked example proposes a
    $\text{top-}k$ retrieval that returns the $k$ nearest stored
    transitions and trains DQN on that subset. My implementation
    instead draws a full mini-batch \emph{stochastically} according to
    the softmax weights, with replacement
    (\texttt{np.random.choice(p=weights)}). This keeps the mini-batch
    size identical to plain DQN and lets the same $\beta$ knob
    smoothly interpolate between near-uniform and near-greedy sampling,
    rather than having a separate discrete $k$.
    \item \textbf{Capacity range.} The worked example suggests
    $N \in \{50, 200, 500\}$. I sweep over $\{1\,000, 10\,000, 100\,000\}$
    because these are the realistic regimes for an online RL replay
    buffer, and so that the largest setting matches the
    Stable-Baselines3 default buffer size.
    \item \textbf{Added RL-side sweep.} The worked example only varies
    bio-inspired and structural knobs. I additionally sweep the DQN
    learning rate, which turned out to interact non-trivially with the
    biased Hopfield sampling and was useful for the discussion.
    \item \textbf{Diagnostic and follow-up.} I track the Hopfield
    energy of the most recent observation throughout training (used in
    Q6 of the discussion), and I validate the best sweep configuration
    with a full 5-seed re-run instead of stopping at the sweep's 3
    seeds.
\end{itemize}

\section{Results}
\label{sec:results}

\subsection{Main experiment}

Figure~\ref{fig:main} and Table~\ref{tab:main} compare the plain DQN
baseline with the default hybrid ($\beta=5$, capacity = buffer size)
over $5$ seeds each.

\begin{figure}[t]
    \centering
    \includegraphics[width=0.92\linewidth]{figures/main_hybrid_vs_baseline.png}
    \caption{Learning curves (20-episode rolling mean, mean $\pm$ std across
    5 seeds) for baseline DQN and the default hybrid on CartPole-v1.}
    \label{fig:main}
\end{figure}

\begin{table}[t]
    \centering
    \begin{tabular}{lcc}
        \toprule
        \textbf{Mode} & \textbf{Final return} & \textbf{Steps to return=400} \\
        \midrule
        Baseline  & $236.4 \pm 98.6$ & $41\,116$ \\
        Hybrid (default, $\beta=5$) & $102.8 \pm 56.4$ & never \\
        \bottomrule
    \end{tabular}
    \caption{Convergence score (mean return over the final 10\% of
    episodes) and sample efficiency (first timestep at which the smoothed
    return hits 400) for the main experiment. Relative improvement of
    hybrid over baseline: $-56.5\%$. Welch $t$-test on final returns:
    $t=-2.35$, $p=0.054$.}
    \label{tab:main}
\end{table}

The default hybrid is clearly worse than the baseline: it never reaches
a smoothed return of $400$, its final return is about $56\%$ lower, and
the difference is borderline significant ($p \approx 0.05$). Three
hyperparameter sweeps are reported next to identify the root cause.

\subsection{Sweep 1 -- Hopfield retrieval temperature $\beta$}

I vary $\beta \in \{1.0, 5.0, 10.0\}$ with $3$ seeds each. Results are in
Figure~\ref{fig:sweep_beta} and the top of Table~\ref{tab:sweeps}.

\begin{figure}[t]
    \centering
    \includegraphics[width=0.92\linewidth]{figures/sweep_beta.png}
    \caption{Effect of the Hopfield retrieval temperature $\beta$ on
    the learning curve (mean $\pm$ std over 3 seeds per value).}
    \label{fig:sweep_beta}
\end{figure}

Lower $\beta$ values produce substantially better learning. $\beta = 1.0$
reaches a smoothed return of about $380$ in the final rolling window,
while $\beta = 10.0$ stays below $100$. The intuition is that at low
$\beta$ the attention weights are almost uniform, so the Hopfield
sampling degrades gracefully toward random replay and preserves DQN's
stability; at high $\beta$ the sampling concentrates on a handful of
nearby patterns, starving the Bellman update of the decorrelated
experience it needs.

\subsection{Sweep 2 -- Hopfield capacity $N$}

I vary the number of patterns the memory can hold, $N \in \{1\,000,\ 10\,000,\ 100\,000\}$
(default). Figure~\ref{fig:sweep_cap} and Table~\ref{tab:sweeps} show
that capacity has no clear effect: the three curves are within each
other's standard-deviation bands throughout training, and the best final
return (at $N=1\,000$) is not much higher than the worst (at $N=10\,000$).
This is a useful negative result: it rules out ``not enough storage'' as
the reason the default hybrid loses.

\begin{figure}[t]
    \centering
    \includegraphics[width=0.92\linewidth]{figures/sweep_capacity.png}
    \caption{Effect of the Hopfield capacity $N$ (mean $\pm$ std over 3
    seeds per value).}
    \label{fig:sweep_cap}
\end{figure}

\subsection{Sensitivity to bio-inspired hyperparameters}

The two bio-inspired knobs are $\beta$ and the Hopfield capacity $N$.
Sweeps 1 and 2 vary them in isolation; the project specification also
asks for a joint sensitivity heatmap so that any interaction between
the two can be inspected. Figure~\ref{fig:sensitivity} shows the
convergence score over a $2 \times 2$ grid of $(\beta, N)$ values
covering the most informative corners of the design space:
$\beta \in \{1.0, 5.0\}$ crossed with $N \in \{1\,000, 100\,000\}$.
Each cell is the mean (and standard deviation) over 3 seeds.

\begin{figure}[t]
    \centering
    \includegraphics[width=0.7\linewidth]{figures/sensitivity_heatmap.png}
    \caption{Sensitivity heatmap: mean final return as a function of
    Hopfield retrieval temperature $\beta$ and capacity $N$, 3 seeds
    per cell. Lighter colours mark higher returns.}
    \label{fig:sensitivity}
\end{figure}

The heatmap is more informative than the one-dimensional sweeps
suggested. The $\beta = 1$ row is no longer uniformly best: at the
small capacity ($N=1\,000$) the score collapses to about $87$, well
below the $\beta = 5$ row at the same capacity, while at the large
capacity ($N=100\,000$) $\beta = 1$ is much better than $\beta = 5$.
In other words, there is a clear \emph{interaction} between the two
hyperparameters. A plausible explanation is that low $\beta$ produces
a soft, near-uniform attention distribution; with only $1\,000$ stored
patterns the resulting mini-batch is drawn from too small a pool to
recover the decorrelation that DQN needs, so the soft averaging hurts.
With $100\,000$ patterns the same soft attention has enough diversity
to behave like a slightly biased random sample, which works well. At
high $\beta$ the attention is already concentrated on a few patterns,
so adding more candidates makes little difference. This explains why
Sweep 2 (which fixed $\beta=5$) showed no effect of capacity: the
effect only appears when $\beta$ is also low.

\subsection{Sweep 3 -- RL learning rate}

I vary $\mathrm{lr} \in \{5\!\times\!10^{-4},\ 2.3\!\times\!10^{-3}\text{ (default)},\ 5\!\times\!10^{-3}\}$.
See Figure~\ref{fig:sweep_lr} and Table~\ref{tab:sweeps}.

\begin{figure}[t]
    \centering
    \includegraphics[width=0.92\linewidth]{figures/sweep_lr.png}
    \caption{Effect of the DQN learning rate on the hybrid's training curve
    (mean $\pm$ std over 3 seeds per value).}
    \label{fig:sweep_lr}
\end{figure}

The hybrid benefits from a smaller learning rate: $5\!\times\!10^{-4}$
reaches about $180$, while the Zoo's default of $2.3\!\times\!10^{-3}$
(optimised for plain DQN) only reaches $110$. This is consistent with
the picture that the biased Hopfield sampling amplifies optimisation
noise, so a more conservative step size helps.

\begin{table}[t]
    \centering
    \small
    \begin{tabular}{llcc}
        \toprule
        \textbf{Sweep} & \textbf{Value} & \textbf{Mean final} & \textbf{Std} \\
        \midrule
        $\beta$       & $1.0$    & $256.8$ & $80.9$ \\
        $\beta$       & $5.0$    & $110.0$ & $71.4$ \\
        $\beta$       & $10.0$   & $66.1$  & $38.1$ \\
        \midrule
        Capacity $N$  & $1\,000$    & $140.5$ & $103.5$ \\
        Capacity $N$  & $10\,000$   & $45.0$  & $23.6$ \\
        Capacity $N$  & $100\,000$  & $110.0$ & $71.4$ \\
        \midrule
        lr            & $5\!\times\!10^{-4}$   & $150.4$ & $24.3$ \\
        lr            & $2.3\!\times\!10^{-3}$ & $110.0$ & $71.4$ \\
        lr            & $5\!\times\!10^{-3}$   & $109.7$ & $20.9$ \\
        \bottomrule
    \end{tabular}
    \caption{All sweeps in one table. ``Mean final'' is the average over
    the last 10\% of episodes, averaged over 3 seeds.}
    \label{tab:sweeps}
\end{table}

\subsection{Best configuration vs.\ baseline}

The sweeps point to $\beta=1.0$, a small-to-default learning rate, and
essentially any capacity. I ran the combination ($\beta=1.0$, default
capacity and learning rate) with the full set of 5 seeds and compared it
to the baseline from the main experiment. Results are in
Figure~\ref{fig:best} and Table~\ref{tab:best}.

\begin{figure}[t]
    \centering
    \includegraphics[width=0.92\linewidth]{figures/best_vs_baseline.png}
    \caption{Best hybrid configuration ($\beta=1.0$) vs.\ plain DQN
    baseline, 5 seeds each, mean $\pm$ std.}
    \label{fig:best}
\end{figure}

\begin{table}[t]
    \centering
    \begin{tabular}{lc}
        \toprule
        \textbf{Mode} & \textbf{Final return} \\
        \midrule
        Baseline              & $236.4 \pm 98.6$ \\
        Best hybrid ($\beta=1.0$) & $263.7 \pm 82.0$ \\
        \midrule
        Relative improvement   & $+11.5\%$ \\
        Welch $t$-test         & $t=0.43$, $p=0.68$ \\
        \bottomrule
    \end{tabular}
    \caption{Final return of the best hybrid configuration compared with
    the baseline. The improvement is positive but not statistically
    significant.}
    \label{tab:best}
\end{table}

With the correct $\beta$, the hybrid's mean final return is about
$11.5\%$ higher than the baseline's. However, the standard-deviation
bands still overlap, and Welch's $t$-test gives $p=0.68$, i.e.\ the
difference is well within seed-to-seed noise. This is an honest result
rather than a disappointing one: it simply says that on CartPole-v1 the
Hopfield buffer is \emph{comparable} to random replay when tuned
correctly, and worse when tuned badly.

\subsection{Computational overhead}

Averaged over all runs, a baseline training takes about $142$ seconds on
a Colab CPU; a hybrid training takes about $184$ seconds. The overhead
is therefore roughly $30\%$, almost all of it spent in the
\texttt{softmax}($\beta X q$) call inside the Hopfield memory. This is a
meaningful cost for a marginal gain.

\section{Discussion}
\label{sec:discussion}

The project specification lists ten evaluation questions. I address each
of them in turn.

\paragraph{Q1: Does the hybrid outperform the baseline in sample efficiency,
and is the difference significant?}
With default hyperparameters, no --- the hybrid is worse and the
difference is borderline significant ($p=0.054$). With tuned
hyperparameters, the hybrid's final return is about $11.5\%$ higher, but
the difference is not statistically significant ($p=0.68$). Neither
configuration reaches the baseline's sample-efficiency target ($400$ in
smoothed return) faster.

\paragraph{Q2: Which bio-inspired hyperparameter has the greatest effect?}
$\beta$ is the dominant knob in absolute terms: moving from $\beta=1$
to $\beta=10$ at the default capacity drops the mean final return from
$257$ to $66$, a roughly fourfold change. Capacity, varied alone at
$\beta=5$, has no monotone effect (Sweep~2). However, the joint
heatmap in Section~4.4 reveals an \emph{interaction}: at low $\beta$
the capacity becomes important, with $N=1\,000$ collapsing the
performance to $87$ while $N=100\,000$ delivers $257$. So the most
accurate one-line answer is that $\beta$ is the primary driver, and
capacity only matters when $\beta$ is small. This is also why the
1-D Sweep 2 (which used $\beta=5$ throughout) found no capacity
effect.

\paragraph{Q3: Does the system converge, or does it drift after an initial peak?}
Looking at Figure~\ref{fig:main}, the baseline exhibits the classical
DQN pattern of rising sharply, hitting a peak around timestep $40\,000$,
and then partially regressing --- a mild form of catastrophic
forgetting. The default hybrid never reaches a peak worth discussing.
The best hybrid (Figure~\ref{fig:best}) continues rising through to the
end of training, suggesting it is slightly less prone to the late-training
regression, though at a cost in peak height.

\paragraph{Q4: Seed sensitivity.}
Standard deviations of the final return are in the $40$--$100$ range,
which is substantial relative to the means. DQN is known to be
seed-sensitive on CartPole (my first single-seed baseline run, reported
in the progress logs, happened to be particularly bad), and the hybrid
inherits this sensitivity. I conservatively used 5 seeds for the main
experiment and 3 for the sweeps.

\paragraph{Q5: What is the best configuration?}
$\beta=1.0$, Hopfield capacity equal to the replay buffer size
($100\,000$), learning rate $2.3\!\times\!10^{-3}$, network
$[256, 256]$, batch size $64$, $\gamma=0.99$, $\mathrm{train\_freq}=256$,
$\mathrm{gradient\_steps}=128$, target update interval $10$, exploration
fraction $0.16$. See
Table~\ref{tab:best} for its performance.

\paragraph{Q6: Do the internal dynamics evolve meaningfully?}
Yes. I recorded the Hopfield energy of \texttt{self.\_last\_obs} at every
50-th \texttt{sample()} call during a representative run. The energy
decreases over training (from roughly $-2$ at the start to roughly $-5$
to $-8$ in the middle) as the memory fills with related patterns, and it
never becomes flat or diverges. This confirms that the memory is doing
nontrivial work rather than returning constants. The trajectory is also
noisy rather than monotone: episodes in which the cart visits a
previously-unseen region of the state space produce temporarily higher
energies, and the subsequent storage of those observations pulls the
energy back down. In effect, the memory's energy landscape is shaped by
the agent's own exploration, which is one of the clearest signs that
the coupling is working as intended. A flat energy trace would suggest
the memory had converged and stopped adapting; a monotonically rising
one would indicate storage saturation. Neither pathology is observed
during the $50\,000$-step runs reported here.

\paragraph{Q7: Is the coupling bidirectional?}
Information flows in both directions. The environment produces
transitions that go into the memory (forward), and the memory dictates
which transitions the DQN optimiser sees (backward). There is no explicit
feedback in the Hopfield update itself --- the memory only stores --- but
the DQN policy's changing behaviour changes what transitions end up in
the memory, which is an indirect form of feedback. Removing this
feedback loop (e.g.\ by pre-collecting a fixed set of transitions and
freezing the memory) would remove the only source of adaptation and is
therefore incompatible with a valid coupling.

\paragraph{Q8: Computational overhead.}
The hybrid is about $30\%$ slower per run. For CartPole this is
negligible in absolute terms. For a large-scale RL problem with a
replay buffer of $10^6$ or $10^7$ transitions and a state dimension in
the thousands, the $O(Nd)$ cost of the Hopfield $\beta X q$ would become
prohibitive. The standard remedies (top-$k$ retrieval, approximate
nearest-neighbour indexing, clustering of stored patterns) would bring
the cost back down and are straightforward engineering extensions.

\paragraph{Q9: Could the Hopfield buffer be replaced by a simpler engineered
substitute?}
On this particular problem, yes --- a uniform random buffer does a
little better than the default Hopfield and only a little worse than the
tuned Hopfield, with no hyperparameters and a tenth of the compute.
A $k$-nearest-neighbour index with a fixed $k$ would very likely match
or beat our best configuration while being more interpretable. What the
Hopfield framing does give us is (i) a single scalar $\beta$ that
interpolates between random and nearest-neighbour sampling, and (ii) an
energy function that provides a diagnostic signal. Those are the
principled features that an engineered substitute would have to
reproduce.

\paragraph{Q10: Would I recommend the hybrid approach in practice?}
On CartPole-v1, \emph{no}. The best hybrid is indistinguishable from
random replay at the statistical-significance level, and it is $30\%$
more expensive. The hybrid would only be worth the extra complexity in
a problem where the relevance of past transitions varies strongly across
the state space (e.g.\ very sparse rewards, large state spaces with
distinct ``regions'', or non-stationary environments). I would tentatively
expect it to pay off on MiniGrid-style environments with sparse rewards
or on multi-task settings; that is an obvious direction for future work.
A second practical caveat is that the Hopfield buffer adds a
hyperparameter ($\beta$) to which DQN turned out to be extremely
sensitive. In a realistic workflow this means any deployment of the
hybrid needs a tuning budget of its own, whereas plain replay has no
such knob. The compute savings of being tuning-free should not be
underestimated. My overall recommendation, based on this study, is to
treat Hopfield retrieval not as a drop-in replacement for random replay
but as an \emph{optional} component to be considered when the problem
structure suggests it will help.

\section{Conclusion}
\label{sec:conclusion}

I built a Hopfield-based replay buffer for DQN, integrated it cleanly
with Stable-Baselines3, and ran a paired comparison against plain DQN on
CartPole-v1 with full hyperparameter sweeps. The main empirical findings
are: (i) default Hopfield retrieval ($\beta=5$) hurts DQN significantly
because it biases the mini-batch and undermines the decorrelation that
uniform replay provides; (ii) retrieval temperature is by far the most
important knob in the system, far more so than capacity or the RL
learning rate; and (iii) when tuned to $\beta=1$, the hybrid matches and
slightly exceeds the baseline, but the difference is not statistically
significant on CartPole. Bio-inspired replay is therefore not a free win
in this setting. The natural next step is to test the same
integration on an environment where random replay is known to be a weak
baseline --- typically one with sparse or deceptive rewards.

\bibliographystyle{plainnat}
\bibliography{references}

\end{document}
